\documentclass[11pt]{article}

% Change "review" to "final" to generate the final (sometimes called camera-ready) version.
% Change to "preprint" to generate a non-anonymous version with page numbers.
\usepackage[preprint]{acl}

% Standard package includes
\usepackage{times}
\usepackage{latexsym}

% For proper rendering and hyphenation of words containing Latin characters (including in bib files)
\usepackage[T1]{fontenc}
% For Vietnamese characters
% \usepackage[T5]{fontenc}
% See https://www.latex-project.org/help/documentation/encguide.pdf for other character sets

% This assumes your files are encoded as UTF8
\usepackage[utf8]{inputenc}

% This is not strictly necessary, and may be commented out,
% but it will improve the layout of the manuscript,
% and will typically save some space.
\usepackage{microtype}

% This is also not strictly necessary, and may be commented out.
% However, it will improve the aesthetics of text in
% the typewriter font.
\usepackage{inconsolata}

%Including images in your LaTeX document requires adding
%additional package(s)
\usepackage{graphicx}
\hbadness=10000
\vbadness=10000

\makeatletter
\renewcommand\outauthor{%
  \begin{tabular}[t]{c}%
  \ifacl@anonymize
    \bfseries Anonymous ACL submission
  \else
    \mdseries\@author
  \fi
  \end{tabular}}
\makeatother

% If the title and author information does not fit in the area allocated, uncomment the following
%
%\setlength\titlebox{<dim>}
%
% and set <dim> to something 5cm or larger.

\title{Evaluating Post-Fine-Tuning Defense Mechanisms Against Harmful Fine-Tuning Attacks}

% Author information can be set in various styles:
% For several authors from the same institution:
% \author{Author 1 \and ... \and Author n \\
%         Address line \\ ... \\ Address line}
% if the names do not fit well on one line use
%         Author 1 \\ {\bf Author 2} \\ ... \\ {\bf Author n} \\
% For authors from different institutions:
% \author{Author 1 \\ Address line \\  ... \\ Address line
%         \And  ... \And
%         Author n \\ Address line \\ ... \\ Address line}
% To start a separate ``row'' of authors use \AND, as in
% \author{Author 1 \\ Address line \\  ... \\ Address line
%         \AND
%         Author 2 \\ Address line \\ ... \\ Address line \And
%         Author 3 \\ Address line \\ ... \\ Address line}

\author{
Neil Pendyala, Anda He, Aniketh Kancherla, Fabien Lazes, Liangbin Zhao \and Yurun Zi \\
University of Wisconsin--Madison \\
\texttt{\{npendyala, ahe28, akancherla, lazes, lzhao258, yzi2\}@wisc.edu}
}

%\author{
%  \textbf{First Author\textsuperscript{1}},
%  \textbf{Second Author\textsuperscript{1,2}},
%  \textbf{Third T. Author\textsuperscript{1}},
%  \textbf{Fourth Author\textsuperscript{1}},
%\\
%  \textbf{Fifth Author\textsuperscript{1,2}},
%  \textbf{Sixth Author\textsuperscript{1}},
%  \textbf{Seventh Author\textsuperscript{1}},
%  \textbf{Eighth Author \textsuperscript{1,2,3,4}},
%\\
%  \textbf{Ninth Author\textsuperscript{1}},
%  \textbf{Tenth Author\textsuperscript{1}},
%  \textbf{Eleventh E. Author\textsuperscript{1,2,3,4,5}},
%  \textbf{Twelfth Author\textsuperscript{1}},
%\\
%  \textbf{Thirteenth Author\textsuperscript{3}},
%  \textbf{Fourteenth F. Author\textsuperscript{2,4}},
%  \textbf{Fifteenth Author\textsuperscript{1}},
%  \textbf{Sixteenth Author\textsuperscript{1}},
%\\
%  \textbf{Seventeenth S. Author\textsuperscript{4,5}},
%  \textbf{Eighteenth Author\textsuperscript{3,4}},
%  \textbf{Nineteenth N. Author\textsuperscript{2,5}},
%  \textbf{Twentieth Author\textsuperscript{1}}
%\\
%\\
%  \textsuperscript{1}Affiliation 1,
%  \textsuperscript{2}Affiliation 2,
%  \textsuperscript{3}Affiliation 3,
%  \textsuperscript{4}Affiliation 4,
%  \textsuperscript{5}Affiliation 5
%\\
%  \small{
%    \textbf{Correspondence:} \href{mailto:email@domain}{email@domain}
%  }
%}

\begin{document}
\maketitle
\begin{abstract}
Large language models (LLMs) are typically released with safety alignment mechanisms that reduce the likelihood of generating harmful outputs. However, recent work shows that these safeguards can be compromised through post-training fine-tuning, even when the fine-tuning dataset contains only a small amount of harmful content. This vulnerability creates a realistic security risk because fine-tuning is widely used to customize models for downstream tasks. In this project, we study the robustness of safety-aligned language models under harmful fine-tuning attacks and evaluate multiple defense mechanisms designed to mitigate this problem. Our experimental pipeline reproduces the alignment stage from the Antidote framework and introduces behavior drift using parameter-efficient LoRA fine-tuning. We then evaluate several defense strategies that operate at different stages of the training pipeline, including alignment-stage defenses (Vaccine, RepNoise), fine-tuning-stage defenses (LISA), and post-fine-tuning defenses (Antidote and Panacea). Safety performance is measured using the BeaverTails dataset, while downstream task utility is evaluated on SST-2, GSM8K, and AGNews. The goal of this work is to systematically compare these defenses within a unified experimental framework and analyze their ability to recover safety behavior while preserving model utility.
\end{abstract}

\section{Introduction}

Large language models (LLMs) are increasingly deployed in real-world applications where safety and reliability are critical. To mitigate risks associated with harmful content generation, many LLMs undergo safety alignment before release. These alignment procedures typically involve supervised fine-tuning or reinforcement learning from human feedback (RLHF), which encourages the model to refuse harmful instructions while maintaining useful task performance.

Despite these safeguards, recent research shows that safety alignment is fragile under post-training fine-tuning. When users fine-tune an aligned model on a new dataset, the model may gradually forget its safety alignment and begin producing unsafe outputs. This phenomenon, often referred to as harmful fine-tuning, can occur even when the harmful portion of the dataset is relatively small. The risk is particularly concerning in fine-tuning-as-a-service scenarios, where model providers allow users to upload custom datasets for training.

To address this issue, several recent works propose defense mechanisms designed to prevent or mitigate harmful fine-tuning. These defenses operate at different stages of the training pipeline, including the alignment stage, the fine-tuning stage, and the post-fine-tuning stage. However, there is currently limited experimental comparison of these approaches within a unified framework.

In this project we implement and evaluate multiple defense strategies against harmful fine-tuning attacks. Our pipeline reproduces the alignment stage using the Antidote framework and introduces behavioral drift through parameter-efficient LoRA fine-tuning. We then compare several representative defense methods that operate at different stages of the training process. Our goal is to analyze how effectively these defenses restore safety behavior while preserving downstream task performance.

\section{Related Work}

Safety alignment aims to ensure that large language models (LLMs) generate responses that follow human preferences for helpfulness and harmlessness. 

The BeaverTails dataset \citep{ji2023beavertails} provides a large-scale benchmark for evaluating safety behavior in LLMs, containing over 330,000 annotated prompts designed to measure harmful and safe responses.


However, several studies have demonstrated that safety alignment can be weakened through harmful fine-tuning. 
Huang et al. \citep{huang2024antidote} show that even a small amount of harmful data introduced during fine-tuning can significantly degrade model safety behavior.

Several defense strategies have been proposed to mitigate harmful fine-tuning attacks. 
Representation Noising (RepNoise) \citep{rosati2024repnoise} introduces noise into internal representations during alignment training in order to weaken the relationship between harmful prompts and model activations. 
LISA \citep{huang2024lisa} proposes a fine-tuning stage defense that constrains parameter drift using proximal regularization while alternating between alignment and user datasets.

Post-fine-tuning defenses attempt to repair models after harmful fine-tuning has already occurred. 
Antidote \citep{huang2024antidote} identifies harmful parameters using Wanda importance scores and removes them using structured pruning. 
Panacea \citep{wang2025panacea} instead learns a perturbation vector applied to the model after fine-tuning that suppresses harmful outputs while preserving task performance.

In addition to safety evaluation, several benchmark datasets are commonly used to evaluate downstream task performance. 
The Stanford Sentiment Treebank (SST-2) \citep{socher2013sst} measures sentiment classification accuracy, while GSM8K \citep{cobbe2021gsm8k} evaluates mathematical reasoning ability through multi-step word problems. 
These datasets allow researchers to measure whether safety defenses degrade the model's general capabilities.

Our work builds on these studies by implementing and comparing multiple defense mechanisms across different stages of the training pipeline under a unified experimental framework. 

\section{Method}

Our experimental pipeline consists of three stages: safety alignment, harmful fine-tuning attack, and defense evaluation as shown in \autoref{fig:pipeline}.

\begin{figure*}[t]
\centering
\includegraphics[width=\textwidth]{experiment_pipeline.pdf}
\caption{Experimental pipeline for evaluating defenses against harmful fine-tuning attacks.}
\label{fig:pipeline}
\end{figure*}

\textbf{Stage 1: Safety Alignment}

We first reproduce the safety alignment training pipeline provided by the Antidote repository. 
The base model used in our experiments is LLaMA-2-7B. 
Alignment training is performed using a subset of the BeaverTails dataset containing safety-filtered prompts and responses.

Training is executed using the provided Antidote script \texttt{sft.sh}, which performs supervised fine-tuning with LoRA adapters. 
The output of this stage is an aligned LoRA checkpoint which serves as the baseline model for subsequent experiments.

\textbf{Stage 2: Harmful Fine-Tuning Attack}

To simulate post-training attacks we perform LoRA fine-tuning under three different attack scenarios inspired by the Antidote evaluation setup.

\begin{itemize}
\item \textbf{Benign forgetting attack}: the model is fine-tuned on benign instruction-following data such as Alpaca. Although the dataset is not harmful, the additional training can cause the model to forget previously learned safety alignment.
\item \textbf{Harmful fine-tuning attack}: the model is fine-tuned on harmful prompts derived from adversarial datasets such as AdvBench, which encourages the model to produce unsafe responses.
\item \textbf{Mixed fine-tuning attack}: the training dataset contains a mixture of benign instruction-following data and harmful prompts, simulating realistic fine-tuning scenarios where harmful examples may appear within otherwise benign datasets.
\end{itemize}

During this stage the base model parameters remain frozen while only the LoRA adapter parameters are updated, allowing us to efficiently simulate behavioral drift while maintaining computational feasibility.

\textbf{Stage 3: Defense Mechanisms}

After harmful fine-tuning we apply multiple defense strategies representing different stages of the training pipeline:

\begin{itemize}
\item Vaccine: alignment-stage defense
\item RepNoise: representation noising defense during alignment
\item LISA: fine-tuning stage defense using proximal regularization
\item Antidote: post-fine-tuning defense that computes Wanda importance scores to identify parameters responsible for harmful behaviors and prunes them to restore safety alignment.
\item Panacea: post-fine-tuning perturbation defense
\end{itemize}

Each defense method will be evaluated under identical experimental conditions.

\textbf{Evaluation}

We measure both safety and utility metrics:

\begin{itemize}
\item Harmful Score (HS): proportion of harmful outputs generated
\item Fine-Tune Accuracy (FA): downstream task accuracy on SST-2, GSM8K, and AGNews
\end{itemize}

\section{Datasets and Experimental Setup}

Safety behavior is evaluated using the BeaverTails dataset \citep{ji2023beavertails}, which contains unique prompts annotated for harmful and safe responses. 
To simulate adversarial attack scenarios we additionally use the AdvBench dataset, which contains adversarial prompts designed to elicit harmful or unsafe responses from language models.

Instruction-following fine-tuning is performed using the Alpaca dataset \citep{taori2023alpaca}, which contains approximately 52,000 instruction-response pairs.

Downstream task performance is evaluated using three additional datasets:

\begin{itemize}
\item SST-2 \citep{socher2013sst} for sentiment classification
\item GSM8K \citep{cobbe2021gsm8k} for mathematical reasoning
\item AGNews \citep{zhang2015charcnn} for news topic classification
\end{itemize}

Safety performance is measured using the Harmful Score (HS), defined as the proportion of model outputs classified as harmful when evaluated on BeaverTails prompts.

Utility performance is measured using top-1 accuracy on the downstream benchmarks.

All experiments will be conducted using a single NVIDIA H100 GPU with 80GB VRAM. 
Alignment training requires approximately 67GB of GPU memory and takes approximately 1.5 hours to complete.

Each experimental run consists of harmful fine-tuning (1–2 hours), pruning or defense application (<10 minutes), and evaluation (~20 minutes). 
The full experimental sweep is estimated to require approximately 20 hours of compute time on a single GPU.

\section{Exploratory Analysis}

Initial work focused on reproducing the Antidote training environment and verifying the alignment training pipeline.

The Antidote repository was cloned and configured, and the alignment training script \texttt{sft.sh} was successfully executed. 
This produced a safety-aligned LoRA checkpoint based on the LLaMA-2-7B base model.

To validate the harmful fine-tuning attack pipeline we conducted preliminary LoRA fine-tuning experiments using the TinyLlama-1.1B model. 
Only the adapter parameters were trained while the base model weights remained frozen.

Preliminary comparisons between the base model and fine-tuned model outputs show measurable behavioral drift, confirming that LoRA fine-tuning can influence model behavior despite modifying a small subset of parameters.

These experiments confirm that the alignment and attack pipelines are functioning correctly and provide the foundation for the full experimental evaluation.

\section{Planned Experiments}

Our planned experiments will compare multiple defense strategies against harmful fine-tuning attacks.

For each defense method we will run a one-factor hyperparameter sweep consisting of:

\begin{itemize}
\item three learning rates
\item three epoch values
\item three harmful dataset ratios
\end{itemize}

Each configuration will be evaluated under three different attack scenarios: benign forgetting, harmful fine-tuning, and mixed fine-tuning. This allows us to analyze how defense mechanisms behave under different forms of safety degradation.

This results in nine experimental configurations per defense method.

Each run will consist of the following pipeline:

\begin{enumerate}
\item Start from the aligned LoRA checkpoint
\item Perform harmful fine-tuning using LoRA adapters
\item Apply the defense mechanism
\item Evaluate safety and utility metrics
\end{enumerate}

Results will be analyzed to determine which defense strategies most effectively restore safety behavior while minimizing degradation of downstream task performance.

\section{Project Timeline}

The remaining work for this project will be completed in four phases.

Phase 1 consists of completing the implementation of the Antidote pruning pipeline and integrating it into our experimental framework.

Phase 2 involves reproducing additional defense mechanisms including Vaccine, RepNoise, LISA, and Panacea.

Phase 3 consists of running the full experimental evaluation across all hyperparameter configurations and collecting safety and utility metrics.

Phase 4 will focus on analyzing experimental results and preparing the final report summarizing the effectiveness of each defense strategy.

\subsection{Task Allocation}

Work on this project is divided among team members based on experimental components.

Neil Pendyala assists with integrating Vaccine into the alignment stage and ensuring compatibility with the harmful fine-tuning attack pipeline. Report organization and document preparation are also coordinated by Neil Pendyala.

Yurun Zi focuses on reproducing the Antidote alignment training environment and generating aligned checkpoints.

Aniketh Kancherla works on implementing LoRA fine-tuning experiments and analyzing behavioral drift in model outputs.

Liangbin Zhao assists with integrating RepNoise into the alignment stage and ensuring compatibility with the harmful fine-tuning attack pipeline.

Fabien Lazes focuses on dataset preparation and evaluation metrics, including computing harmful scores and downstream benchmark accuracy.

Anda He assists with experimental evaluation and result analysis across defense methods.

% Bibliography entries for the entire Anthology, followed by custom entries
%\bibliography{custom,anthology-overleaf-1,anthology-overleaf-2}

% Custom bibliography entries only

\bibliography{custom}

\end{document}